\documentclass[11pt,a4paper]{article}

\title{\textbf{Ten Analytical Business Questions}\\[0.5cm] \textbf{AeroGrid}\\[0.3cm] Low-Altitude Airspace Intelligence Warehouse}
\author{AeroGrid Development Team}
\date{\today}

\begin{document}

\maketitle

\section{Introduction}
To ensure the AeroGrid is built to answer real operational and strategic planning needs rather than generic reporting, its dimensional model is driven by ten concrete, measurable analytical questions. Each question directly supports low-altitude traffic management, safety oversight, and infrastructure planning.

\section{The Ten Business Questions}

\begin{enumerate}
    \item \textbf{Altitude Layer \& Grid Crowding:} \\
    How crowded do different altitude layers and map grids become throughout a typical day and night in selected urban or airport areas? \\
    \textit{Significance:} Helps traffic managers identify temporal and spatial congestion hotspots to prevent airspace bottlenecks.

    \item \textbf{Longest-Occupied Sky Zones:} \\
    Which parts of the low-altitude sky stay occupied the longest and therefore need the most careful capacity management? \\
    \textit{Significance:} Highlights permanent high-traffic corridors requiring dedicated routing rules or automated monitoring.

    \item \textbf{Traffic Volume \& Height Trends:} \\
    How has the overall volume of low-altitude traffic and the average flying height changed from month to month and across seasons? \\
    \textit{Significance:} Essential for tracking industry growth and adapting infrastructure to shifting operational altitudes.

    \item \textbf{Peak Times \& Traffic Mix Shifts:} \\
    At what times of day or on which days of the week does traffic peak, and does the mix of manned aircraft versus drones shift noticeably? \\
    \textit{Significance:} Informs dynamic staffing and separation standards for mixed-traffic airspace environments.

    \item \textbf{Close Encounter Frequency \& Hotspots:} \\
    How many close encounters occur each day, and where in the sky (by height and location) are they most common? \\
    \textit{Significance:} Directly measures safety performance and pinpoints hazardous airspace sectors requiring intervention.

    \item \textbf{Risk-Prone Vehicle Categories:} \\
    Which types of vehicles appear most often in close encounters, and are certain operators or aircraft categories involved more frequently? \\
    \textit{Significance:} Guides targeted safety training, regulatory policy updates, and collision-avoidance technology mandates.

    \item \textbf{Impact of Flight Restrictions (TFRs):} \\
    When a temporary flight restriction is active, how much does traffic density drop or shift in the surrounding area? \\
    \textit{Significance:} Measures compliance levels and evaluates how effectively traffic reroutes around restricted zones.

    \item \textbf{Weather Correlation with Risk and Operations:} \\
    Do stronger winds, lower visibility, or lower cloud ceilings reduce the number of operations or increase the number of close encounters? \\
    \textit{Significance:} Crucial for establishing weather-based operational limits and advisories for unmanned flights.

    \item \textbf{Corridor Capacity Stress-Testing:} \\
    If the number of drones in a busy corridor were suddenly doubled or tripled, how would density and the number of close encounters change? \\
    \textit{Significance:} Provides predictive scenario analysis for future traffic scaling without compromising safety.

    \item \textbf{Candidate Identification for New Drone Corridors:} \\
    Which altitude bands and geographic zones currently show both low traffic and few historical close encounters, making them the strongest candidates for new recommended drone corridors? \\
    \textit{Significance:} Enables proactive infrastructure planning to safely expand commercial drone delivery and transit routes.
\end{enumerate}

\end{document}
